\documentclass[11pt]{article}
\usepackage[T1]{fontenc}
\usepackage{lmodern,amsmath,amssymb,amsthm,microtype}
\usepackage[margin=1in]{geometry}
\usepackage{booktabs}
\usepackage[hidelinks]{hyperref}
\hypersetup{pdftitle={F-purity of the maximal permanental ideal of a generic three-by-four matrix},pdfauthor={Ying Xie},pdfcreator={},pdfproducer={}}
\ifdefined\pdfinfoomitdate\pdfinfoomitdate=1\fi
\ifdefined\pdfsuppressptexinfo\pdfsuppressptexinfo=15\fi
\ifdefined\pdftrailerid\pdftrailerid{}\fi
\newtheorem{theorem}{Theorem}[section]
\newtheorem{proposition}[theorem]{Proposition}
\newtheorem{lemma}[theorem]{Lemma}
\newtheorem{corollary}[theorem]{Corollary}
\theoremstyle{remark}\newtheorem{remark}[theorem]{Remark}
\DeclareMathOperator{\per}{per}
\DeclareMathOperator{\rank}{rank}
\DeclareMathOperator{\htop}{ht}
\DeclareMathOperator{\CT}{CT}
\newcommand{\F}{\mathbb F}
\newcommand{\A}{\mathbb A}
\newcommand{\Z}{\mathbb Z}
\newcommand{\m}{\mathfrak m}
\numberwithin{equation}{section}
\title{$F$-purity of the maximal permanental ideal\\of a generic three-by-four matrix}
\author{Ying Xie\\Kennesaw State University\\\texttt{yxie2@kennesaw.edu}}
\date{}
\begin{document}
\maketitle
\begin{abstract}
Let $X$ be a generic $3\times4$ matrix over a field of odd
characteristic $p$, and let $I$ be the ideal generated by its four
maximal permanents. We prove that $k[X]/I$ is $F$-pure if and only if
$p\equiv1\pmod6$, resolving a conjecture of Chau. The proof computes
the coefficient that occurs in Fedder's criterion. Eliminating one row
reduces this coefficient to a quartic determinant, whose constant-term
description identifies it with a truncated Domb-number sum. A finite
binomial transformation gives the exact value
$\binom{(p-1)/2}{(p-1)/6}^{2}$ when $p\equiv1\pmod6$, and zero
otherwise. We include an elementary complete-intersection argument
in every odd characteristic and a telescoping proof of the classical
finite transformation used in the calculation.
\end{abstract}

\section{Introduction}
For a matrix $X=(x_{ij})$ of independent variables, write $P_t(X)$ for
the ideal generated by the permanents of its $t\times t$ submatrices.
In odd characteristic these ideals behave differently from their
determinantal counterparts. Even the ideals of maximal permanents of
rectangular matrices present difficult questions about height,
reducedness, and Frobenius singularities. Boralevi, Carlini, Micha\l{}ek,
and Ventura study their codimensions in \cite{BCMV}; Chau studies
their $F$-singularities in \cite{Chau}.

Consider the generic matrix
\[
 X=\begin{pmatrix}
 a_1&a_2&a_3&a_4\\ b_1&b_2&b_3&b_4\\v_1&v_2&v_3&v_4
 \end{pmatrix}.
\]
Let $f_j$ be the permanent of the $3\times3$ matrix obtained by
deleting column $j$, and set $I=(f_1,f_2,f_3,f_4)$. Chau's
Conjecture~4.5 in the preprint version of \cite{Chau} predicts the
following classification, which was checked there through
characteristic $31$.

\begin{theorem}\label{thm:main}
Let $k$ be any field of characteristic $p>2$. For the generic
$3\times4$ matrix $X$, the ring $k[X]/P_3(X)$ is $F$-pure if and only
if $p\equiv1\pmod6$.
\end{theorem}

Here $F$-purity means purity of the Frobenius homomorphism: tensoring
the map $R\longrightarrow F_*R$, $a\longmapsto F_*(a^p)$, with any
$R$-module preserves injectivity. The $R$-module structure on $F_*R$
is $a\cdot F_*b=F_*(a^pb)$.

The central calculation is more precise than the nonvanishing
statement needed for Theorem~\ref{thm:main}. Put
\[
 H=f_1f_2f_3f_4,\qquad
 M=\prod_{j=1}^{4}a_jb_jv_j,\qquad
 c(p)=[M^{p-1}]H^{p-1}\pmod p.
\]
Brackets denote coefficient extraction in the indicated polynomial.

\begin{theorem}\label{thm:coefficient}
For every odd prime $p$,
\begin{equation}\label{eq:maincoefficient}
 c(p)=
 \begin{cases}
 \displaystyle\binom{(p-1)/2}{(p-1)/6}^{\!2}&p\equiv1\pmod6,\\[6pt]
 0&p\not\equiv1\pmod6
 \end{cases}
 \quad\text{in }\F_p.
\end{equation}
\end{theorem}

The proof identifies $c(p)$ with
\begin{equation}\label{eq:sumintro}
 \sum_{n=0}^{p-1}\frac{D_n}{4^n},\qquad
 D_n=\sum_{j=0}^{n}\binom nj^2\binom{2j}{j}
                         \binom{2n-2j}{n-j}.
\end{equation}
The integers $D_n$ are the Domb numbers. Their transformations and
congruences are established number theory. In particular, Mao and Liu
prove stronger congruences modulo $p^3$ in \cite{ML}. The contribution
here is the reduction of the permanental coefficient to
\eqref{eq:sumintro}. We evaluate the sum modulo $p$ using the finite
transformation of Sun \cite[Lemma~3.1]{Sun}, also recorded as equation~(2.5) in
the preprint version of \cite{ML}. Appendix~\ref{app:telescoping}
supplies a proof of that transformation, so no supercongruence is
needed as an input.

All references to numbered statements of Chau and Mao--Liu use the
specific preprint versions linked in the bibliography. The corresponding
journal publication data are also supplied.

\section{The complete-intersection and coefficient criteria}
\subsection{An elementary height computation}
Let $A$ consist of the first two rows of $X$. Expansion along the last
row gives
\begin{equation}\label{eq:L}
 (f_1,f_2,f_3,f_4)^{\mathsf t}=L(A)v,\qquad
 L_{ii}=0,\quad L_{ij}=a_kb_l+a_lb_k\quad(i\ne j),
\end{equation}
where $\{k,l\}=\{1,2,3,4\}\setminus\{i,j\}$ and
$v=(v_1,v_2,v_3,v_4)^{\mathsf t}$. The matrix $L$ is symmetric.

\begin{proposition}\label{prop:ci}
Over every field of characteristic different from two, the four
polynomials $f_1,f_2,f_3,f_4$ form a regular sequence. In particular,
$I$ has height four.
\end{proposition}
\begin{proof}
It suffices to compute dimensions over an algebraic closure.
For $0\le r\le4$, let $S_r$ be the locus in the eight-dimensional
$A$-space on which $L$ has rank $r$. We show
\begin{equation}\label{eq:rankbound}
 \dim S_r\le 4+r.
\end{equation}
For $r=4$ this is immediate. The determinant of $L$ is a nonzero
polynomial, as the specialization below shows, so the rank-at-most-three
locus has codimension at least one.

The rank-at-most-two locus is contained in the common zero set of
$\det L$ and the principal minor $2L_{12}L_{13}L_{23}$. These two
polynomials have no common irreducible factor. Each $L_{ij}$ is an
irreducible polynomial of the form $xy+zw$ in four independent
variables. To see that $L_{12}$ does not divide $\det L$, specialize
the four columns of $A$ to
\[
 (1,0)^{\mathsf t},\quad (0,1)^{\mathsf t},\quad
 (1,1)^{\mathsf t},\quad (1,-1)^{\mathsf t}.
\]
Then
\[
 L_{12}=0,\quad(L_{13},L_{14},L_{23},L_{24},L_{34})=(1,1,-1,1,1),
\]
and
\[
 \det L=(L_{13}L_{24}-L_{14}L_{23})^2=4\ne0.
\]
Permuting columns gives the same conclusion for the other factors.
Since height-one prime ideals in a polynomial unique factorization
domain are principal, the indicated common zero set has codimension
at least two. This proves \eqref{eq:rankbound} for $r=2$.

A symmetric matrix with zero diagonal and rank at most one is zero:
its principal $2\times2$ minors are $-L_{ij}^2$. It remains to bound
the dimension of $L=0$. These equations say that all $2\times2$
permanents of $A$ vanish. If some column is $(a,b)^{\mathsf t}$ with
$ab\ne0$, every other column has the form $t_j(a,-b)^{\mathsf t}$.
The permanent of two such columns is $-2ab\,t_jt_k$; hence at most one
of the other columns is nonzero. Each of these finitely many loci has
dimension at most three. Otherwise every column is on a coordinate
axis. Nonzero columns on opposite axes have nonzero permanent, so all
nonzero columns must lie on the same axis. These loci have dimension
four. Thus $\dim S_0\le4$, and the rank-one stratum is empty.

Above $S_r$, equation \eqref{eq:L} has a vector-space fiber of
dimension $4-r$. Consequently every rank stratum in $V(I)$ has
dimension at most $(4+r)+(4-r)=8$. The height theorem for four
generators gives the reverse bound in the twelve-dimensional ambient
space. Thus $\htop I=4$. A list of four homogeneous generators of
height four in a polynomial ring is a regular sequence, since the
ring is Cohen--Macaulay.
\end{proof}

\subsection{Fedder's criterion}
For an ideal $J$ in characteristic $p$, write $J^{[p]}$ for the ideal
generated by $p$th powers of its elements. We use the following form
of Fedder's criterion \cite{Fedder}: if $S$ is a localization of a
polynomial ring over $\F_p$ at a maximal ideal $\m$, and $I\subseteq\m$,
then $S/I$ is
$F$-pure if and only if $(I^{[p]}:I)\not\subseteq\m^{[p]}$.
For a regular sequence $f_1,\ldots,f_s$, the colon ideal is
\begin{equation}\label{eq:colon}
 (I^{[p]}:I)=I^{[p]}+\bigl((f_1\cdots f_s)^{p-1}\bigr).
\end{equation}
We apply this first over $\F_p$ at the homogeneous maximal ideal.

\begin{proposition}\label{prop:fedder}
For every field $k$ of odd characteristic $p$, the ring $k[X]/I$ is
$F$-pure if and only if $c(p)\ne0$.
\end{proposition}
\begin{proof}
The polynomial $H^{p-1}$ has total degree $12(p-1)$. A monomial
outside the ideal generated by the $p$th powers of the twelve
variables must have every exponent at most $p-1$. Its degree
therefore forces it to be $M^{p-1}$. Proposition~\ref{prop:ci} and
\eqref{eq:colon} show that $c(p)=0$ implies failure of $F$-purity at
the origin over $\F_p$.

For sufficiency, there is a global splitting. Choose an additive map
$\lambda:k\to k$ satisfying
$\lambda(a^pb)=a\lambda(b)$ and $\lambda(1)=1$.
Such a map exists by choosing a $k^p$-basis of $k$ containing $1$,
projecting to its $1$-coordinate, and taking a $p$th root.
On a monomial in the twelve variables define
\[
 T(bX^{\alpha})=
 \begin{cases}
 \lambda(b)X^{(\alpha-(p-1)\boldsymbol1)/p},
       &\alpha_i\equiv p-1\pmod p\text{ for every }i,\\
 0,&\text{otherwise}.
 \end{cases}
\]
Then $T(a^pg)=aT(g)$ and $T(H^{p-1})=c(p)$. If $c(p)\ne0$,
the map $\phi(g)=c(p)^{-1}T(H^{p-1}g)$ satisfies $\phi(1)=1$.
Moreover,
\[
 H^{p-1}f_jg=f_j^p\prod_{i\ne j}f_i^{p-1}g,
\]
so $\phi$ preserves every $(f_j)$ and hence $I$. It induces a
Frobenius splitting of $k[X]/I$, which implies $F$-purity.

Finally, $F$-purity descends along faithfully flat ring extensions.
Thus $F$-purity over $k$ would imply $F$-purity over $\F_p$, since
all the equations are defined over $\F_p$. The necessity already
proved over $\F_p$ therefore applies to every $k$.
\end{proof}

\section{Eliminating one row}
We now work over $\F_p$. The elementary power sums are
\[
 \sum_{a\in\F_p}a^e=
 \begin{cases}
 -1,&e>0\text{ and }p-1\mid e,\\
 0,&\text{otherwise},
 \end{cases}
\]
where the case $e=0$ is the sum of $p$ copies of $1$.
Consequently, for a homogeneous polynomial $G$ of degree $N(p-1)$,
\begin{equation}\label{eq:powersum}
 \sum_{x\in\F_p^N}G(x)=(-1)^N
 [x_1^{p-1}\cdots x_N^{p-1}]G.
\end{equation}
Indeed, a monomial can survive only if every exponent is a positive
multiple of $p-1$, and total degree then forces equality in every
coordinate.

\begin{proposition}\label{prop:row}
With $L$ as in \eqref{eq:L},
\begin{equation}\label{eq:rowcoefficient}
 c(p)=\left[\prod_{j=1}^4a_j^{p-1}b_j^{p-1}\right]
                       \det L(A)^{p-1}.
\end{equation}
\end{proposition}
\begin{proof}
Apply \eqref{eq:powersum} in twelve variables and sum over $v$ first.
If $L(A)$ is singular, each nonempty fiber of the map
$v\mapsto L(A)v$ has cardinality divisible by $p$. The sum of
$\prod_j(L(A)v)_j^{p-1}$ over all $v$ is therefore zero in $\F_p$.
If $L(A)$ is invertible, a change of variables gives
\[
 \sum_{v\in\F_p^4}\prod_{j=1}^4(L(A)v)_j^{p-1}
 =\prod_{j=1}^4\sum_{w\in\F_p}w^{p-1}=1.
\]
In either case the sum is $\det L(A)^{p-1}$. Applying
\eqref{eq:powersum} in the remaining eight variables proves
\eqref{eq:rowcoefficient}.
\end{proof}

\subsection{The quartic determinant}
Make the formal Laurent-polynomial substitution $a_i=b_it_i$ and
write $e_j=e_j(t_1,t_2,t_3,t_4)$ for the elementary symmetric
polynomials. Since each variable column has degree two in $\det L$,
\begin{equation}\label{eq:detquartic}
 \det L(A)=\left(\prod_{i=1}^4b_i^2\right)D(t),\qquad
 D(t)=-4(e_1e_3-4e_4).
\end{equation}
Here is an explicit verification. For a symmetric matrix with zero
diagonal and upper-triangular entries $(a,b,c,d,e,f)$ in row order,
the determinant is
\[
 a^2f^2+b^2e^2+c^2d^2-2(abef+acdf+bcde).
\]
After removing the $b_i$ factors from $L$, the six entries are
\[
 t_3+t_4,\ t_2+t_4,\ t_2+t_3,\
 t_1+t_4,\ t_1+t_3,\ t_1+t_2.
\]
Substitution in the displayed determinant gives
$-4(e_1e_3-4e_4)$. Equivalently, its twelve monomials of exponent
pattern $(2,1,1,0)$ have coefficient $-4$, and all other coefficients
are zero.

If $t_i$ has exponent $\alpha_i$ in $D(t)^{p-1}$, then after
substitution the exponents on $a_i,b_i$ are respectively
$\alpha_i,2(p-1)-\alpha_i$. Thus
\begin{equation}\label{eq:quarticextract}
 c(p)=[(t_1t_2t_3t_4)^{p-1}]D(t)^{p-1}.
\end{equation}
This is a polynomial coefficient calculation; it does not discard
points at which some $b_i$ vanishes.

Set
\[
 Q(t)=\left(\sum_{i=1}^4t_i\right)
      \left(\sum_{i=1}^4t_i^{-1}\right).
\]
Since $e_1e_3=e_4Q$ and $(-4)^{p-1}=1$ in $\F_p$,
\eqref{eq:quarticextract} becomes
\begin{equation}\label{eq:ct}
 c(p)=\CT(Q-4)^{p-1}
     =\sum_{n=0}^{p-1}4^{-n}\CT Q^n.
\end{equation}
The second equality follows from
$\binom{p-1}{n}\equiv(-1)^n\pmod p$.

\begin{lemma}\label{lem:domb}
For every nonnegative integer $n$, $\CT Q^n=D_n$.
\end{lemma}
\begin{proof}
The multinomial theorem gives
\[
 \CT Q^n=\sum_{n_1+n_2+n_3+n_4=n}
             \binom{n}{n_1,n_2,n_3,n_4}^{\!2}.
\]
Group terms according to $j=n_1+n_2$, and use
$\sum_{a=0}^j\binom ja^2=\binom{2j}{j}$ for each of the two pairs.
The resulting sum is precisely the definition of $D_n$ in
\eqref{eq:sumintro}.
\end{proof}

\section{Evaluation of the Domb-number sum}
We use the integer identity
\begin{equation}\label{eq:transform}
 D_n=\sum_{k=0}^{\lfloor n/2\rfloor}
 \binom{n+k}{3k}\binom{2k}{k}^{\!2}\binom{3k}{k}4^{n-2k}.
\end{equation}
It is a classical transformation; see \cite[Lemma~3.1]{Sun} and
\cite[equation~(2.5)]{ML}. A complete telescoping proof is given in
Appendix~\ref{app:telescoping}.

\begin{proposition}\label{prop:sum}
For an odd prime $p$, let $C_p=\sum_{n=0}^{p-1}D_n/4^n\in\F_p$.
If $p\not\equiv1\pmod3$, then $C_p=0$. If $p\equiv1\pmod3$ and
$m=(p-1)/3$, then
\begin{equation}\label{eq:binomial}
 C_p=16^{-m}\binom{2m}{m}^{\!2}\ne0.
\end{equation}
\end{proposition}
\begin{proof}
Insert \eqref{eq:transform}, interchange the finite sums, and use
the elementary summation identity
$\sum_{n=2k}^{p-1}\binom{n+k}{3k}=\binom{p+k}{3k+1}$.
This gives, first as an equality of rational numbers and then modulo
$p$,
\begin{equation}\label{eq:finiteouter}
 C_p=\sum_{k=0}^{(p-1)/2}16^{-k}\binom{2k}{k}^{\!2}
                  \binom{3k}{k}\binom{p+k}{3k+1}.
\end{equation}
The product of the last two binomial coefficients is
\begin{equation}\label{eq:valuation}
 \binom{3k}{k}\binom{p+k}{3k+1}
 =\frac{\prod_{j=-2k}^{k}(p+j)}{k!(2k)!(3k+1)}.
\end{equation}
Because $0\le k\le(p-1)/2$, the numerator has exactly one factor
divisible by $p$, namely the factor with $j=0$; it is exactly $p$.
The factors $k!$ and $(2k)!$ are units modulo $p$. Also
$1\le3k+1<2p$, so the only possible cancellation of that factor
is $3k+1=p$.

If $p\not\equiv1\pmod3$, there is no such $k$ and every summand
vanishes. This includes $p=3$. Otherwise put $m=(p-1)/3$. At
$k=m$ we cancel $p$ in \eqref{eq:valuation} over the integers before
reducing; its value modulo $p$ is
\[
 \frac{(-1)^{2m}(2m)!m!}{m!(2m)!}=1.
\]
Only this summand remains, proving \eqref{eq:binomial}. The binomial
coefficient is nonzero modulo $p$ because $2m<p$.
\end{proof}

\begin{proof}[Proof of Theorems~\ref{thm:main} and \ref{thm:coefficient}]
Equation~\eqref{eq:ct} and Lemma~\ref{lem:domb} give $c(p)=C_p$.
For odd $p$, the condition $p\equiv1\pmod3$ is equivalent to
$p\equiv1\pmod6$. Proposition~\ref{prop:fedder} and
Proposition~\ref{prop:sum} prove Theorem~\ref{thm:main}.

For the exact coefficient, use the rational identity
\[
 \binom{2m}{m}=(-4)^m\binom{-1/2}{m}.
\]
When $m=(p-1)/3$, all denominators are invertible modulo $p$.
Reducing $-1/2$ to $(p-1)/2$ and using binomial symmetry gives
\[
 16^{-m}\binom{2m}{m}^{\!2}
 =\binom{(p-1)/2}{m}^{\!2}
 =\binom{(p-1)/2}{(p-1)/6}^{\!2}\quad\text{in }\F_p.
\]
This proves Theorem~\ref{thm:coefficient}.
\end{proof}

\section{Consequences and comparison}
When $p\equiv1\pmod6$, the proof supplies an explicit compatible
Frobenius splitting, and hence reducedness of the maximal permanental
quotient. In every odd characteristic, independently of $F$-purity,
Proposition~\ref{prop:ci} gives the Hilbert series
\[
 \operatorname{Hilb}_{k[X]/I}(z)=\frac{(1-z^3)^4}{(1-z)^{12}}.
\]
Failure of $F$-purity in the other characteristics is not, by itself,
a statement that the quotient is nonreduced.

The residue calculation is consistent with the stronger arithmetic
result of Mao--Liu \cite[Theorem~1.1]{ML}. For $p>3$ that theorem
implies, modulo $p$, that the sum in \eqref{eq:sumintro} equals
$4x^2$ if $p=x^2+3y^2\equiv1\pmod3$, and zero if
$p\equiv2\pmod3$. Thus it provides a second way to conclude the
nonvanishing criterion once the permanental reduction is established.
The proof above uses only the finite transformation and does not
assert a new congruence for the Domb numbers.

For illustration, the first few coefficient values are
\[
\begin{array}{c|rrrrrrrr}
p&3&5&7&11&13&17&19&23\\\hline
c(p)&0&0&2&0&4&0&7&0.
\end{array}
\]
Exact finite checks of the determinant identity and the telescoping
identities accompany the manuscript. The classification for all odd
primes follows from the proofs, rather than from extrapolation of
these values.

\appendix
\section{A telescoping proof of the finite transformation}
\label{app:telescoping}
This appendix proves the established identity \eqref{eq:transform}
over the integers. The proof is included to make the arithmetic input
directly checkable.

Put
\begin{align*}
 f(n,k)&=\binom nk^2\binom{2k}{k}\binom{2n-2k}{n-k},
       &&0\le k\le n,\\
 g(n,k)&=\binom{n+k}{3k}\binom{2k}{k}^{\!2}
                       \binom{3k}{k}4^{n-2k},
       &&0\le2k\le n,
\end{align*}
and extend each summand by zero outside its stated range. For $n\ge2$
define
\begin{align*}
 B_n&=2(2n-1)\bigl(5n(n-1)+2\bigr),\\
 P(n,k)&=4k^3+(6-18n)k^2+(26n^2-15n)k+10n^2-12n^3,\\
 r(n,k)&=\frac{k^3P(n,k)}{n^2(2n-2k-1)},\\
 s(n,k)&=\frac{-12(3n-2)k^4}{(n+k)(n+k-1)}.
\end{align*}
The two telescoping identities are
\begin{align}
 &n^3f(n,k)-B_nf(n-1,k)+64(n-1)^3f(n-2,k)\notag\\
 &\hspace{25mm}=r(n,k+1)f(n,k+1)-r(n,k)f(n,k),\label{eq:tel1}\\
 &n^3g(n,k)-B_ng(n-1,k)+64(n-1)^3g(n-2,k)\notag\\
 &\hspace{25mm}=s(n,k+1)g(n,k+1)-s(n,k)g(n,k).\label{eq:tel2}
\end{align}
We give polynomial forms of these identities to avoid an implicit
appeal to symbolic summation software. In \eqref{eq:tel1} put $h=n-k$.
Cancellation of factorials and clearing denominators reduces it to
\begin{align}
 &2(2k+1)h^3P(n,k+1)-2k^3(2h-3)P(n,k)\notag\\
 &=2n^5(2h-1)(2h-3)-B_nh^3(2h-3)
                         +32(n-1)h^3(h-1)^3.\label{eq:poly1}
\end{align}
In \eqref{eq:tel2} put $h=n-2k$ and $C=-12(3n-2)$. The analogous
polynomial identity is
\begin{align}
 &C\bigl((n+k-1)(2k+1)h(h-1)-8k^4\bigr)\notag\\
 &=8n^3(n+k)(n+k-1)-2B_nh(n+k-1)
                         +32(n-1)^3h(h-1).\label{eq:poly2}
\end{align}
Both \eqref{eq:poly1} and \eqref{eq:poly2} follow by expansion in
$\Z[n,k]$. To recover the summand identities, the required successive
ratios are
\begin{align*}
 \frac{f(n,k+1)}{f(n,k)}
  &=\frac{(2k+1)(n-k)^3}{(k+1)^3(2n-2k-1)},\\
 \frac{f(n-1,k)}{f(n,k)}
  &=\frac{(n-k)^3}{2n^2(2n-2k-1)},\\
 \frac{f(n-2,k)}{f(n,k)}
  &=\frac{(n-k)^3(n-k-1)^3}
 {4n^2(n-1)^2(2n-2k-1)(2n-2k-3)},\\
 \frac{g(n,k+1)}{g(n,k)}
  &=\frac{(n+k+1)(2k+1)(n-2k)(n-2k-1)}{8(k+1)^4},\\
 \frac{g(n-1,k)}{g(n,k)}&=\frac{n-2k}{4(n+k)},\\
 \frac{g(n-2,k)}{g(n,k)}
  &=\frac{(n-2k)(n-2k-1)}{16(n+k)(n+k-1)}.
\end{align*}
The factors in the numerators account for zero summands at the upper
boundary. None of the displayed denominators in the $f$ ratios
vanishes on its integer summation range when $n\ge2$; the same is
true of the $g$ ratios. The certificates are finite at the endpoints,
and $r(n,0)=s(n,0)=0$. At the upper endpoint plus one, the extended
summand is zero. Summing \eqref{eq:tel1} or \eqref{eq:tel2} therefore
gives the recurrence
\[
 n^3U_n=B_nU_{n-1}-64(n-1)^3U_{n-2}\qquad(n\ge2).
\]
Both finite sums have $U_0=1$ and $U_1=4$. Over $\mathbb Q$, these
initial values and the nonzero coefficient $n^3$ determine the
sequence uniquely. The sums agree as integers for every $n$, proving
\eqref{eq:transform}. Reduction modulo a prime is performed only
after this integer identity has been established.

\paragraph{Acknowledgment.}
The author acknowledges OpenAI GPT-6 for assistance with mathematical
exploration, exact computations, proof organization, and manuscript
preparation.

\begingroup\small
\begin{thebibliography}{5}\raggedright
\bibitem{BCMV}
A.~Boralevi, E.~Carlini, M.~Micha\l{}ek, and E.~Ventura,
\emph{On the codimension of permanental varieties},
Advances in Mathematics \textbf{461} (2025), article 110079.
\href{https://doi.org/10.1016/j.aim.2024.110079}{doi:10.1016/j.aim.2024.110079}.

\bibitem{Chau}
T.~Chau, \emph{The $F$-singularities of algebras defined by permanents},
Collectanea Mathematica \textbf{77} (2026), 843--859.
\href{https://doi.org/10.1007/s13348-025-00494-8}{doi:10.1007/s13348-025-00494-8}.
Preprint: \href{https://arxiv.org/abs/2509.09980v1}{arXiv:2509.09980v1}.

\bibitem{Fedder}
R.~Fedder, \emph{$F$-purity and rational singularity},
Transactions of the American Mathematical Society \textbf{278}
(1983), no.~2, 461--480.
\href{https://doi.org/10.1090/S0002-9947-1983-0701505-0}{doi:10.1090/S0002-9947-1983-0701505-0}.

\bibitem{ML}
G.-S.~Mao and Y.~Liu, \emph{Proof of some conjectural congruences
involving Domb numbers and binary quadratic forms},
Journal of Mathematical Analysis and Applications \textbf{516}
(2022), no.~1, article 126493.
\href{https://doi.org/10.1016/j.jmaa.2022.126493}{doi:10.1016/j.jmaa.2022.126493}.
Preprint under the title \emph{Proof of some conjectural congruences
involving Domb numbers}:
\href{https://arxiv.org/abs/2112.00511v3}{arXiv:2112.00511v3}.

\bibitem{Sun}
Z.-H.~Sun, \emph{Congruences for Domb and Almkvist--Zudilin numbers},
Integral Transforms and Special Functions \textbf{26} (2015),
no.~8, 642--659.
\href{https://doi.org/10.1080/10652469.2015.1034122}{doi:10.1080/10652469.2015.1034122}.
\end{thebibliography}
\endgroup
\end{document}
